\section{Threats to Validity}

Construct validity is limited by the current label source. The artifact set has deterministic baseline outputs for field typing, silver-v2 development annotations, and strict overlays from blind Codex reviewers, but no human gold. The five discrepancy categories are useful for analysis, yet their boundaries can be subjective, especially between representation\_discrepancy and incomplete for affected\_versions. The older URL-only LLM draft is not treated as an evaluation label set, and the dual-Codex overlays are explicitly marked \texttt{label\_is\_human=false}.

The routing contract itself is also a construct-validity assumption. The method maps equivalent and representation differences to normalization, incompleteness to augmentation, temporal differences to timeline interpretation, and factual conflicts to evidence-constrained adjudication. That mapping is methodologically useful, but it is not yet validated against practitioner decisions or downstream scanner behavior. Until such validation exists, the routing contract should be read as an auditable baseline design, not as a proven operational policy.

Internal validity is limited by evidence quality. The evidence cache stores fetched text snippets and metadata, not a full reconstruction of dynamic pages. Some URLs return 403, 404, or timeout, so the evidence set is incomplete by construction. The evidence-source diagnostic shows that this incompleteness is source-class dependent: affected\_versions has complete fetched text for NVD pages in the sample, but substantially lower fetched-text rates for GitHub commit/repository pages and GitHub advisories. These failures can bias adjudication toward sources that are easier to fetch. The silver-v2 labels also depend on the fetched evidence context; if important vendor or commit evidence is missing, both the silver label and the baseline comparison can be biased.

Reviewer independence is not evidence independence. In the nine v2 strict FC-source rows, both reviewers cite the same single URL in five cases, and only two cases contain primary/ecosystem evidence cited by both reviewers. Three rows collectively rely only on NVD record evidence. The strict overlay may therefore favor evidence that is shared or easier to fetch, even when the reviewer processes are isolated.

The affected\_versions holdouts reduce prediction-to-label leakage but do not remove all dependence. V1 excludes the development CVEs; v2 excludes both earlier 100-CVE cohorts. Reviewer inputs omit prior labels and sealed predictions, yet all rows are sampled from the same deterministic factual-conflict candidate miner. Each reviewer pair also shares the same Codex model family, prompt contract, and evidence snapshot, so their errors are not independent. In these non-human diagnostics, v1 label/source kappa is 0.268/0.392 with 0.35 strict joint coverage. V2 label/artifact kappa is 0.535/0.669, but strict type coverage is 0.41 and strict FC-source coverage over the full cohort is only 0.09. Reported performance is therefore conditional on selective, easier-to-agree subsets.

External validity is limited by coverage. The adjudication diagnostics cover severity with 80 silver samples, affected\_versions with a 100-row development sample, and affected\_versions with two additional 100-row CVE-disjoint holdouts. Both holdouts remain conditional on the current NVD-GHSA corpus and candidate miner. The affected\_versions predictors are not semantic version-range adjudicators, and affected-version identification itself is known to depend on package, range, and ecosystem semantics \cite{OSVSchema2026,Chen2025AffectedVersions}. Neither field has gold-backed metrics. The current numbers should not be generalized to all fields or vulnerability databases without the pending human audit and cross-field experiments.

Conclusion validity is limited by the selective setup. RQ2 currently reports diagnostic rule-trigger counts rather than performance metrics. Accuracy and macro-F1 are reported against non-human labels, and coverage matters because both adjudicator disagreement and method abstention are selective. The v1 factual-conflict/non-conflict split was discovered after unsealing, so its 16-row factual-conflict analysis is post-hoc. V2 preregisters separate endpoints, but its type method predicts only three strict rows and its source endpoint contains only nine strict rows. V1 and v2 are both now unsealed and must not be used to tune a replacement method.

The RQ2 CWE audit seals predictions before reviews and covers the complete 17-row candidate impact set, but it is not a representative sample of the 8,066 aligned pairs. Sixteen CVEs are disjoint from the primary seed, yet the rule family was selected after earlier AI-candidate error inspection, and all reviewers remain Codex runs from the same model family. The evidence-enhanced stage is additionally selected from first-stage disagreements and regressions; only 28/36 reference snapshots were fetched successfully, so evidence availability can affect adjudication. Combining stages raises strict coverage from 11/17 to 15/17, but two rows remain unresolved and four evidence-supported rows favor the current factual-conflict decision. The seed-disjoint bootstrap interval reaches zero and the exact sign diagnostic is 0.1796. These properties permit a bounded candidate-direction statement, not human-gold accuracy, independent confirmation, or a production switch.

The references audit also covers a complete rule-impact set rather than a representative sample. The two reviewers share one model family and evidence snapshot, and the grouped input necessarily reveals which URLs are being tested together even though the candidate transformation and label are masked. Thirty probes return 404, making the resource-identity definition consequential. The first protocol seal is retained but excluded because its validator contained a confidence constraint absent from the prompt; revision 2 uses fresh reviewers and complete sealed role prompts, but it was created after pilot output existed and is not fully preregistered before any model response. The 32-row audited profile is then selected after revision-2 unsealing. Its candidate-label gains are therefore post-hoc development diagnostics, not independent or human-gold performance.

The references human-review packet removes the hidden-contract problem by requiring each reviewer and resolving author to record the resource-identity definition, but it does not itself resolve construct validity. All 56 rows are still pending and no real human has signed one. Its 24-row priority flag identifies the known definition-sensitive subset without supplying either Codex verdict; this is workflow readiness, not human validation.

A separate current-snapshot RQ2 typing-stability cohort excludes all 717 CVEs previously exposed through development, audits, or holdouts and seals six executable profile predictions for 1,250 globally unique CVEs. After those exclusions, all six prediction columns are identical, so the completed review estimates only baseline stability and cannot identify candidate-profile gain. In this non-human diagnostic, the A/B roles are composed from 28 and 67 disjoint ephemeral sessions using the same model/config, not two persistent or human annotators. A later immutable-log audit identifies three excess B requests and 90 retry row-items. Every requested row eventually has an accepted output, but the old log lacks response-error events and rejected JSON; for one duplicate payload, it cannot identify which attempt failed. The new artifact records this missingness without backfilling an error reason. Accepted rows pass the sealed merge contract, but this retry/provenance gap and same-model dependence limit interpretation. Independent comparison of the references and CWE candidates requires a new time cohort collected after profile freeze, and the current consensus still requires human validation.

Both acquisitions are post-seal by collection time but not by bilateral event time. The second acquisition adds 74 NVD rows, including 39 published after the freeze, but adds no reviewed GHSA row and leaves the one-GHSA matched field views byte-identical. No aligned record therefore has both NVD and GHSA publication timestamps later than the profile seal, so the 250-row fallback remains explicitly snapshot-external and no new cohort is frozen. This source asymmetry means that repeated acquisition can audit readiness but cannot by itself establish temporal validation. The fallback's two passes share the same model, prompt, executable, and evidence projection; disjoint sessions and reverse order reduce state and ordering dependence but not model-family dependence. Only three sampled rows distinguish profiles, and two lack strict consensus in the sealed main review. The observed +1 candidate paired count is consequently sparse, selective same-model development evidence. A subsequent evidence-secondary diagnostic selects exactly those three known differences; two fresh same-family passes reach 3/3 strict candidate-direction agreement, but this post-selection result cannot be substituted into the main 250-row metric. V1 and v2 are excluded wholesale after fail-closed checks reject a confidence/uncertainty mismatch and a nonliteral CWE path, so the accepted v3 output avoids silent repair but remains dependent on an interface refined after failure. None of these stages can establish future-snapshot generalization, a temporal effect, human correctness, or a production-switch benefit.

The all-50 CWE evidence audit removes explicit profile-difference flags from reviewer inputs and covers the entire sealed field, but it remains post-hoc because field selection and all three contract versions follow A/B and profile unsealing. Its 134/135 successful URL snapshots are broad but not proof that the retrieved pages contain complete mapping evidence. V1 confounds literal subset with semantic compatibility, and V2 discovers the literal-quotation gap only at merge; both are archived and excluded, while V3 is accepted only after the interface is refined. The 20 sessions are mutually disjoint but still use one model family. Consequently, 49/50 strict coverage and the 45/49 versus 48/49 comparison are useful construct and control diagnostics, not a preregistered accuracy estimate. They cannot replace the sealed 185/231 versus 186/231 result, establish human correctness, or justify production promotion.

The final 16-row non-CWE stage is also outcome-selected after A/B unsealing. Its selector and thresholds are fixed before its own evidence retrieval, its 50 frozen records all fetch successfully, and G/H sessions are disjoint from prior stages, but G/H still share the same model family and prompt lineage. The staged 238/250 candidate preferentially retains already strict rows and adds decisions only from difficult revealed subsets; its 0.952 coverage and 188/238 versus 191/238 profile agreement therefore cannot estimate full-cohort accuracy. The latter difference is inherited from the earlier CWE stage, since the four new non-CWE decisions do not distinguish profiles. Successful retrieval also does not prove completeness of package mappings, range endpoints, CVSS semantics, or resource identity. The fixed resolution gate fails at 4/16, so neither the 12 remaining abstentions nor the no-go may be converted into human truth.

The paired-outcome and paired-test diagnostics are label-free and exhaustive over the frozen five-label universe, but their 125 assignments are not equally likely empirical outcomes. Treating the 40 candidate-better, 40 current-better, and 45 tied cases as probabilities would impose an unsupported uniform prior. The +/-1.2-point bound applies only to full-cohort accuracy. The exact-test result proves only that no labeling of these three prediction differences can reject the paired null at alpha 0.05; it does not establish method equivalence. The future-cohort counts additionally assume independent sampling and a stationary 3/250 prediction-difference rate, while the power table conditions on an assumed candidate-direction probability among correctness-discordant rows. Neither diagnostic bounds macro-F1, category-specific recall, construct validity, the effect of a revised profile, or the true power of a future cohort without those assumptions.

The eligible-universe census removes the unsupported stationary 3/250 extrapolation for this revealed snapshot, but introduces no correctness evidence. Its 34 union differences are deterministic outputs from profiles developed with earlier diagnostics; the universe is snapshot-external rather than strict event-time, and its contents and rule-trigger strata are now revealed. The census cannot be converted into a confirmatory cohort or an unbiased estimate for a later snapshot. As a label-free diagnostic, current-status stratification explains why the selected slice is not self-weighting, but does not by itself explain every within-stratum deviation in a small sample. Most importantly, this diagnostic's 29 CWE prediction differences and five or three references differences are not correctness-discordant counts. A disagreement-enriched future review cannot support absolute performance or prevalence without a design declared before labels, CVE-level grouping, and a separately probability-sampled control layer.

The two complete-difference audits add labels only after the universe and outcomes are revealed. Their completeness prevents within-impact sampling loss but does not make the selected strata representative, event-time eligible, or independent of profile development. The CWE reviewers share a model family and all rows expose the same disjoint taxonomy structure; the small conditional p-value is therefore not a preregistered significance result. The references worklists hide direct source fields but raw URLs can reveal CVEs and the alias class, and complete-partition consensus is sensitive to how related advisory, NVD, and third-party pages are ontologically grouped. Revision 1 also contains a disclosed merge-code failure; revision 2 is valid only as a repaired development protocol created after that failure. Neither result is human gold, and neither may replace the blank full-cohort or full-impact real-person packets.

The reviewer-C tiebreak reduces missing candidate coverage but adds post-selection dependence. Its worklist is blind to A/B content and its sessions are disjoint, yet the 103 rows are selected only after A/B unsealing and all three roles share one model family, prompt, and snapshot. The 66 resolved rows are therefore an opportunistic same-family majority candidate, not a preregistered independent third expert. The fixed gate correctly remains no-go at 0.6408 selected-row resolution and 0.9704 combined coverage. Reporting baseline agreement on the 1,213 selected candidates does not estimate accuracy, and the unresolved 37 rows may be systematically harder than the retained rows.

The D/E evidence-secondary audit is selected after both A/B and C are unsealed and therefore cannot repair that selection dependence. Its URL selector is frozen before new retrieval and both reviewer inputs are blind to prior labels, but all five roles still share one model family and the 37-row cohort is explicitly the hardest remainder. Nine non-human rows receive no successful non-empty evidence record, while successful page retrieval does not imply complete field evidence. Exact non-human D/E agreement on 25/37 is dominated by 19 shared uncertain decisions; only 6 non-human rows are evidence-qualified strict, and non-human affected\_versions resolution remains 0/28. These values characterize abstention and evidence limitations in a selected cohort, not independent-expert agreement, prevalence, or human-gold accuracy.

The residual non-affected audit is even more targeted: it selects the final two CWE rows and one references row after D/E outcomes are known, and its source, taxonomy, and URL profile shapes are inspected before v1. The 1/3 mechanism-supported result therefore cannot estimate residual resolution probability or validate the \texttt{factual\_conflict} candidate. Absence of a local \texttt{try} supports an uncaught-exception path but does not prove that resource consumption is impossible elsewhere. Likewise, a prohibited CWE category does not prove the alternative source assignment, and a malformed-URL repair is a construct choice rather than an observed redirect. All three promoted fields remain null and combined coverage is unchanged.

The staged frontier is also post-hoc and conditions each marginal ratio on the preceding stage's failures. Its 66/103 and 6/37 yields describe different, progressively harder selected sets and must not be read as a controlled model comparison or extrapolated to real reviewers. Token totals cover successful responses only because the three missing B attempts have no usage record. The stop rule is valid only for further same-model escalation on this revealed cohort; it does not establish that future models, real reviewers, or an independent snapshot would have the same yield.

The non-human diagnostic holdout also reveals an evaluator-contract threat. The current severity baseline uses canonical labels, whereas the reviewer protocol defines the field over labels, scores, vectors, and CVSS versions; 165 strict disagreements therefore combine comparator error with a deliberate difference in construct. The affected\_versions projection separately removes package-specific unbounded records whose only range marker is \texttt{introduced=0}, producing 25 strict one-sided-presence disagreements. A post-hoc profile that encodes these observed mechanisms reaches 1,121/1,147 on the same non-human strict rows, but this is in-sample diagnostic fit and cannot estimate generalization. The blank full-cohort human packet reduces selection bias in the planned audit, yet it has 0/1,250 signed rows and its validator cannot prove that supplied identities correspond to independent real people; identity and review independence require external verification.

This non-human diagnostic shows that protocol-generation drift prevents retrospective pooling. Replaying the fresh severity contract on the older AI-gold rows reverses direction, but both generations still use related Codex workflows, so the reversal does not identify which contract is correct. For affected\_versions, the old seed stores empty projected range values for ten one-sided-package rows; current raw aligned data cannot be substituted as the label-time input. Consequently, neither a pooled accuracy estimate nor current-snapshot reconstruction is valid evidence for the mechanism. A shared real-human calibration under one explicit contract must precede any new candidate seal.

The two AI contract calibrations reduce input and ordering variation but do not remove model-family or selection dependence. Both draw from the already reviewed fresh cohort and require prior non-human strict consensus for selection. V2 is row-disjoint from v1, but its refined clauses were derived after v1 unsealing, so it validates reviewer stability under the revised wording rather than external correctness. The conditional XWiki evidence secondary and both deterministic projection passes are selected after v2 unsealing. Projection v2 closes the earlier coordinate gap with official product POMs and source continuity, but the chosen class set, edge schema, and release parser were developed on this known failure. Its \texttt{incomplete} candidate is therefore a traceable case explanation, not an independent estimate, human agreement, or gold accuracy. The subsequent generic graph contract is fixed before cross-case evidence retrieval, but both cross-case sources had already been reviewed and derive from a non-human-label-conditioned calibration. Equality selection makes the first 8/8 result an expected construct-consistency diagnostic. The non-equal selector avoids reviewer files and excludes XWiki, but its 2/5 agreement diagnostic still compares against same-model-family labels rather than human truth. Registry catalogs additionally implement extensional snapshot semantics and can omit prereleases, as Graylog demonstrates. The later Codex contract candidate explicitly chooses that extensional interpretation after the no-go is known and tests only one InLong row whose two component ranges are identical. The three-row unseen-ecosystem cohort improves selection independence because it comes from the full aligned input rather than consensus-conditioned calibration data, but its ecosystems and structural filters were chosen after the earlier no-go, it contains one row per ecosystem, and all initial product projections abstain. The Deno lockfile recovery contract is fixed before its new evidence fetch, yet it is explicitly designed after inspecting that row's dependency-constraint failure and reuses the same revealed CVE. Its 71 complete mappings support a one-case mechanism recovery, not prevalence, accuracy, or an unbiased success rate. Generalization still requires a human-approved range contract, edge-class mapping rules tested beyond one project, a later untouched cohort, and real-person validation.

The 28-row edge audit is also post-unsealing. Its family rules and structural score allocate development effort; they do not estimate a family success probability. Mattermost then undergoes two explicit protocol revisions before v3 set analysis: v1 exposes an API pagination assumption, and v2 shows that supplied version boundaries need not be GitHub Release objects. V3 narrows the construct to 19 exact Git tags after those failures are known. The cache-only verifier establishes the reported manifest and ancestry facts, but the 0/2 result remains a selected two-row mechanism diagnostic. It cannot be interpreted as Mattermost ground truth, a Go-ecosystem failure rate, or an unbiased evaluation of graph projection.

LF Edge EVE is likewise selected after the 28-row audit and undergoes disclosed protocol discovery before v1. The 207-tag grammar, absent root module, nested component manifests, and pseudo-commit changed paths are known when the contract is frozen, so v1 explicitly disables candidate promotion. The Git pack and independent verifier establish the reported ref, object, timestamp, path, and ancestry facts, but the component and patched-anchor gates compare a public advisory API snapshot with detail observed in a different repository-advisory interface. Interface drift or representation differences may explain that mismatch. The 0/2 result is therefore a selected mechanism and evidence-interface diagnostic, not EVE ground truth, a historical-advisory correction, or a Go-repository prevalence estimate.

Hutool is also selected from the revealed 28-row audit, and its equal 214-token catalogs, 209-token stable grammar, and three aggregate anchors are inspected before v1 is frozen. The independent verifier establishes the cached catalog, POM, JAR, and set facts, but source/JAR containment is inspected at three critical anchors rather than all 209 releases. More importantly, the projection explicitly chooses current-snapshot extensional semantics for two GHSA ranges with no upper bound. The 2/2 mechanism pass and two development \texttt{incomplete} candidates therefore cannot be interpreted as corrected labels, candidate coverage, historical advisory intent, Maven-family prevalence, or human agreement. They only justify evaluating the pre-frozen mechanism on a later blind cohort.

The six-row Hutool external application is disjoint by CVE from the 1,967-record prior-exposure union, and neither baseline nor reviewer labels enter selection. Nevertheless, it is not an untouched temporal holdout: the records come from the same aligned snapshot and their availability and structural routes were inspected before sealing. Its 6/6 projection result therefore supports only retrospective mechanism reuse under the declared snapshot-extensional contract. It cannot estimate future-snapshot coverage, Maven-family prevalence, taxonomic accuracy, or human agreement, and none of its six candidates is promoted.

The leave-one-cohort-out analysis also remains post-hoc: its authority classes, structured features, models, and thresholds were selected after v2 unsealing, and all three cohorts use non-human labels from related workflows. CVE-disjoint splits reduce row leakage but do not remove shared candidate-miner, evidence-source, or model-family dependence. The pooled type gain is not stable on held-out Phase D, and the source models do not exceed the strongest named baseline. We use this analysis only as a no-go gate, not as confirmatory model selection.

The v2 reviewer files also required disclosed mechanical repairs by their original reviewers: one JSON null was converted to an empty string, and short source rationales were expanded without changing labels, sources, or evidence. The frozen merge contract's character limits were not stated explicitly in the reviewer prompt. The repairs do not alter the substantive decisions, but the prompt/validator mismatch is a protocol defect and weakens any claim of a fully hands-off annotation run.

Reproducibility depends on preserving generated artifacts and manifests. The paper points to generated JSON, Markdown, CSV, and SVG outputs so counts can be rechecked, but final publication still needs a stable dependency list, final bibliography format, and a clear statement of which data snapshots were used.
